\subsection{商空间与子空间的对偶}

在本节中，E 表示一个局部凸空间，M 表示 E 的一个向量子空间，而 \(M^{\circ}\) 表示 M 在 E 的对偶 \(E'\) 中的正交集。设 p 是从 E 到 E/M 上的典范映射；则 \(^{t}p\) 是单射的，其像为 \(M^{\circ}\)，从而定义一个向量空间同构（非拓扑的）

\[
\pi : (\mathrm{E} / \mathbf {M}) ^ {\prime} \rightarrow \mathbf {M} ^ {\circ}.
\]

类似地，设 i 是从 M 到 E 中的典范单射。则 \({}^{t}i\) 是满射的（II，第 24 页，命题 2）；其核等于 \(M^{\circ}\)，于是我们得到一个向量空间同构（非拓扑的）

\[
\iota : \mathrm{E} ^ {\prime} / \mathrm{M} ^ {\circ} \rightarrow \mathrm{M} ^ {\prime}.
\]

命题 9. --- (i) 为使 (E/M)' 的子集 A 等度连续，必要且充分的是 \(\pi(A)\) 是 \(E'\) 的等度连续子集。

\begin{enumerate}
\def\labelenumi{(\roman{enumi})}
\setcounter{enumi}{1}
\item
  设 \(\mathfrak{S}\) 是 E 的有界子集组成的一个集，\(\mathfrak{S}_1\) 是诸子集 \(\mathrm{A} \in \mathfrak{S}\) 在 E/M 中的像组成的集。则 \(\pi\) 是从 \((\mathrm{E} / \mathrm{M})_{\mathfrak{S}_1}'\) 到 \(\mathbf{M}^\circ\) 上的一个同构，其中 \(\mathbf{M}^\circ\) 赋予由 \(\mathrm{E}_{\mathfrak{S}}'\) 的拓扑诱导的拓扑。
\item
  设 E 是赋范空间，则 \(\pi\) 是从赋范空间 (E/M)' 到 E' 的赋范子空间 M° 上的一个等距。
\end{enumerate}

设 A 是 \((\mathrm{E}/\mathrm{M})'\) 的子集，\(\mathbf{B} = {}^{t}p(\mathbf{A}) \subset \mathbf{E}'\)。令

\[
q (\xi) = \sup _ {\xi^ {\prime} \in \mathbf {A}} | \langle \xi , \xi^ {\prime} \rangle |
\]

对所有 \(\xi \in \mathrm{E} / \mathbf{M}\) 成立。为使 A 等度连续，必要且充分的是：映射 \(q\)（从 \(\mathrm{E} / \mathbf{M}\) 到 \(\overline{\mathbf{R}}_+\) 中）是连续半范数。这蕴含 \(q \circ p\) 是 \(\mathrm{E}\) 上的连续半范数（II，第 27 页，命题 5，(ii)）。由于我们有

\[
(q \circ p) (x) = \sup _ {x ^ {\prime} \in \mathbf {B}} | \langle x, x ^ {\prime} \rangle |
\]

对所有 \(x \in E\) 成立，而这又蕴含 B 在 \(E'\) 中等度连续，于是 (i) 得证。设 \(A \in S\)，f 是 E/M 上的连续线性型。对每个 \(\lambda \in R_{+}\)，\(|f| \leqslant \lambda\) 在 \(p(A)\) 上成立当且仅当 \(|^{t}p(f)| \leqslant \lambda\) 在 A 上成立；由此得 (ii)。

最后证明 (iii)。设 \(y'\) 属于 \((\mathrm{E} / \mathbf{M})'\)。\(\mathrm{E} / \mathbf{M}\) 中一个元素的范数 \(< 1\)，当且仅当它是 \(p\) 下某个元素的像，该元素具有范数 \(< 1\) 且位于 \(\mathbf{E}\) 中。于是

\[
\begin{array}{l} \| y ^ {\prime} \| = \sup _ {y \in \mathrm{E} / \mathbf {M}, \| y \| <   1} \left| \langle y, y ^ {\prime} \rangle \right| = \sup _ {x \in \mathrm{E}, \| x \| <   1} \left| \langle p (x), y ^ {\prime} \rangle \right| \\ = \sup _ {x \in \mathrm{E}, \| x \| <   1} \left| \langle x, ^ {t} p (y ^ {\prime}) \rangle \right| = \left\| ^ {t} p (y ^ {\prime}) \right\|, \end{array}
\]

且 \({}^t p\) 诱导一个从 (E/M)' 到 \({\mathbf{M}}^{ \circ  }\) 上的等距。

命题 10. --- (i) 为使 M' 的子集 A 等度连续，必要且充分的是：它是 E' 的某个等度连续子集在 \({}^t i\) 下的像。

\begin{enumerate}
\def\labelenumi{(\roman{enumi})}
\setcounter{enumi}{1}
\item
  设 M 在 E 中闭。设 S 是由 E 的有界子集组成的一个覆盖，\(S_{1}\) 是 M 中形如 \(M \cap A\)（A 属于 S）的子集组成的集。双射线性映射 \(\iota\)（从 \(E_{S}^{\prime}/M^{\circ}\) 到 \(M_{S_{1}}^{\prime}\) 上）是连续的。若 S 对于关系 \(\subset\) 是有向集，且由对于 \(\sigma(\mathrm{E}, \mathrm{E}')\) 闭、凸且紧的集组成，则它是一个同胚。
\item
  若假定 E 是赋范的，则 \(\iota\) 是从 \(E'/M^{\circ}\) 到 \(M'\) 上的一个等距。
\end{enumerate}

\({}^t i\) 把 \(\mathbf{E}'\) 的等度连续子集映为 \(\mathbf{M}'\) 的等度连续子集（IV，第 7 页，命题 7）。反之，设 \(\mathbf{A}\) 是 \(\mathbf{M}'\) 的等度连续子集。\(\mathbf{M}\) 的拓扑由限制到 \(\mathbf{M}\) 的 \(\mathbf{E}\) 上连续半范数全体所定义。于是存在连续半范数 \(p\)（在 \(\mathbf{E}\) 上），使得 \(|f(x)| \leqslant p(x)\)（对所有 \(f \in \mathbf{A}\) 及所有 \(x \in \mathbf{M}\)）。设 \(\mathbf{B}\) 是所有线性型 \(g\)（在 \(\mathbf{E}\) 上）组成的集，它们满足 \(|g| \leqslant p\) 且其在 \(\mathbf{M}\) 上的限制属于 \(\mathbf{A}\)。集 \(\mathbf{B}\) 在 \(\mathbf{E}'\) 中等度连续；由 Hahn-Banach 定理（II，第 23 页，推论 1），我们有 \({}^t i(\mathbf{B}) = \mathbf{A}\)，由此得 (i)。

现在证明 (ii)。由 IV 第 6 页的命题 6，线性映射 \({}^t i\)（从 \(\mathbf{E}_{\mathfrak{S}}^{\prime}\) 到 \(\mathbf{M}_{\mathfrak{S}_1}^{\prime}\) 中）连续，并且通过过渡到商，定义了连续线性映射 \(\iota\)（从 \(E_{\mathfrak{S}}^{\prime}/M^{\circ}\) 到 \(M_{\mathfrak{S}_{1}}^{\prime}\) 上）。设 T 是 \(M^{\prime}\) 上的这样一个拓扑：它由 \(E_{\mathfrak{S}}^{\prime}/M^{\circ}\) 的拓扑通过 \(\iota\) 转移而来；它比 \(S_{1}\)-拓扑细。

现在设 \(\mathfrak{S}\) 对于 \(\subset\) 是有向集，且由对于 \(\sigma(\mathrm{E},\mathrm{E}')\) 闭、凸、均衡且紧的集组成。为证明 \(\iota\) 是同胚，即 \(\mathcal{T}\) 比 \(\mathfrak{S}_1\)-拓扑（\(\mathbf{M}'\) 上的）粗，只需证明 \(\mathcal{T}\) 与 \(\mathbf{M}'\) 和 \(\mathbf{M}\) 之间的对偶相容，并且 \(\mathbf{M}\) 中（把 \(\mathbf{M}\) 视为带有 \(\mathcal{T}\) 的对偶）的每个等度连续集都包含在属于 \(\mathfrak{S}_1\) 的某个集的同位相似像中。由于 \(\mathcal{T}\) 比 \(\mathfrak{S}_1\)-拓扑细，且 \(\mathfrak{S}_1\) 是 \(\mathbf{M}\) 的一个覆盖，故线性型 \(y' \mapsto \langle y, y' \rangle\)（\(\mathbf{M}'\) 上的）对于 \(\mathcal{T}\) 连续，这里 \(y \in \mathbf{M}\) 任意。设 \(f\) 是 \(\mathbf{M}'\) 上对于 \(\mathcal{T}\) 连续的线性型；则 \(f \circ {}^t i\) 是 \(\mathrm{E}_{\mathfrak{S}}'\) 上的连续线性型。\(\mathfrak{S}\)-拓扑（\(\mathrm{E}'\) 上的）比 Mackey 拓扑 \(\tau(\mathrm{E}',\mathrm{E})\) 粗；因为映射 \(d_{\mathrm{B}}: \mathrm{E} \to \mathrm{E}'^*\) 对于拓扑 \(\sigma(\mathrm{E},\mathrm{E}')\) 与 \(\sigma(\mathrm{E}'^*,\mathrm{E}')\) 连续，且后者是 Hausdorff 的，故 \(d_{\mathrm{B}}\) 把对于 \(\sigma(\mathrm{E},\mathrm{E}')\) 紧的集映为对于 \(\sigma(\mathrm{E}'^*,\mathrm{E}')\) 紧的集。由 IV 第 3 页的引理 1，存在 \(x_0 \in E\)，使得 \(f(^t i(x')) = \langle x_0, x' \rangle\)（对所有 \(x' \in E'\)）。特别地，\(\langle x_0, x' \rangle = 0\)（对所有 \(x' \in M^\circ\)）；由于 \(\mathbf{M}\) 在 \(\mathrm{E}\) 中闭，我们有 \(x_0 \in \mathbf{M}\)（II，第 45 页，推论 2）；最后，\(f(y') = \langle x_0, y' \rangle\)（对所有 \(y' \in M'\)）。这就证明了 \(\mathcal{T}\) 与 \(\mathbf{M}\) 和 \(\mathbf{M}'\) 之间的对偶相容。

现在设 A 是 M 的一个子集，对于 \(M'\) 上的拓扑 T 等度连续。由 T 的定义并注意到 S 有向这一假设，这意味着存在包含 0 的集 \(B \in S\)，使得 \(\lambda\)（诸数 \(\left|\langle y, x'\rangle\right|\) 当 \(y \in A\) 且 \(x' \in B^{\circ}\) 时的上确界）是有限的（III，第 19 页，命题 7）。由于 B 在 E 中闭，双极定理（II，第 44 页，定理 1）表明我们有 \(A \subset \lambda(B \cap M)\)；这就完成了 (ii) 的证明。

现在证明 (iii)。设 \(y' \in M'\)。我们将证明公式

\[
\| y ^ {\prime} \| = \inf _ {{}^t i (x ^ {\prime}) = y ^ {\prime}} \| x ^ {\prime} \|.\tag{6}
\]

由 IV 第 7 页的命题 8 (ii)，我们有 \(\| {}^t i\| = \| i\|\)，于是 \(\| {}^t i\| \leqslant 1\)，且

\[
\| y ^ {\prime} \| \leqslant \inf _ {{}^t i (x ^ {\prime}) = y ^ {\prime}} \| x ^ {\prime} \|.\tag{7}
\]

由 Hahn-Banach 定理（II，第 23 页，推论 3），存在 E 上的线性型 \(x_0'\)，它延拓 \(y'\) 且范数相同；于是我们得到与 (7) 相反的不等式，因为 \({}^t i(x_0') = y'\)。

注. --- 我们知道（II，第 48 页，命题 7，(ii)），\(\iota\) 是从 \(\mathrm{E}_s^{\prime}|\mathbf{M}^{\circ}\) 到 \(\mathbf{M}_s^\prime\) 上的拓扑向量空间同构（弱对偶）。对于紧凸收敛拓扑，命题 10 表明：当 E 是 Hausdorff 的且 M 在 E 中闭时，\(\iota\) 是从 \(\mathrm{E}_{cc}^{\prime}/\mathrm{M}^{\circ}\) 到 \(\mathbf{M}_{cc}^{\prime}\) 上的同构。对于强拓扑，\(\iota\) 是从 \(\mathrm{E}_b^{\prime}/\mathrm{M}^{\circ}\) 到 \(\mathbf{M}_b^{\prime}\) 中的连续映射；若 E 是 Banach 空间 * 或 E 是半自反的且 M 在 E 中闭（IV，第 15 页）*，则它是同构；但当 E 是 Fréchet 空间时未必如此（IV，第 58 页，习题 5，c)）。

命题 11. --- (i) E/M 上的弱化拓扑是 E 上的弱化拓扑的商拓扑；M 上的弱化拓扑由 E 的弱化拓扑诱导。

\begin{enumerate}
\def\labelenumi{(\roman{enumi})}
\setcounter{enumi}{1}
\tightlist
\item
  E/M 上的 Mackey 拓扑是 E 上的 Mackey 拓扑的商拓扑；M 上的 Mackey 拓扑比由 \(\tau(\mathrm{E}, \mathrm{E}')\) 诱导的拓扑细。
\end{enumerate}

断言 (i) 由 II 第 48 页的命题 7 得出。

典范单射 \(i:M\to E\) 对于弱化拓扑连续，从而对于 Mackey 拓扑 \(\tau(\mathbf{M},\mathbf{M}^{\prime})\) 与 \(\tau(\mathbf{E},\mathbf{E}^{\prime})\) 连续（IV，第 7 页，命题 7）。类似地，典范投影 \(p:E\to E/M\) 对于 Mackey 拓扑连续。我们立即看出，由 \(\tau(\mathbf{E},\mathbf{E}^{\prime})\) 得到的 E/M 上的商拓扑与 E/M 和 \((\mathrm{E}/\mathrm{M})'\) 之间的对偶相容，故由 Mackey 定理（IV，第 2 页，定理 1），它比 E/M 上的 Mackey 拓扑粗。这就证明了 (ii)。

